\section{Procedura operativa}

\subsection{Setup iniziale - Clonare le repository}

Per iniziare a lavorare al progetto è necessario clonare entrambe le repository. Eseguire i seguenti comandi:

\subsubsection{Clone della repository di codice}
\begin{verbatim}
git clone https://github.com/eghosa02/mirofish_trading_offline.git
cd mirofish_trading_offline
\end{verbatim}

\subsubsection{Clone della repository di documentazione}
\begin{verbatim}
git clone https://github.com/eghosa02/mirofish_trading_offline_doc.git
cd mirofish_trading_offline_doc
\end{verbatim}

\subsection{Creazione di una nuova feature}

\subsubsection{Step 1: Identificare il numero della issue}
Ogni nuova feature deve essere collegata a una issue su GitHub. Annotare il numero della issue (es. \#42).

\subsubsection{Step 2: Creare il branch locale}
\begin{verbatim}
git checkout main
git pull origin main
git checkout -b feature#<numero_issue>
\end{verbatim}

Esempio per la issue \#42:
\begin{verbatim}
git checkout -b feature#42
\end{verbatim}

\subsubsection{Step 3: Sviluppare la feature}
Effettuare le modifiche necessarie sul branch creato. È possibile effettuare commit locali durante lo sviluppo:

\begin{verbatim}
git add .
git commit -m "Descrizione breve delle modifiche"
\end{verbatim}

\subsection{Caricamento dei cambiamenti}

\subsubsection{Step 1: Verificare i cambiamenti}
Prima di eseguire il push, verificare lo stato dei file modificati:

\begin{verbatim}
git status
\end{verbatim}

\subsubsection{Step 2: Aggiungere i file modificati}
\begin{verbatim}
git add .
\end{verbatim}

Se si desiderano aggiungere solo specifici file:
\begin{verbatim}
git add <percorso/file>
\end{verbatim}

\subsubsection{Step 3: Eseguire il commit}
\begin{verbatim}
git commit -m "Messaggio descrittivo del commit"
\end{verbatim}

\subsubsection{Step 4: Eseguire il push}
\begin{verbatim}
git push origin feature#<numero_issue>
\end{verbatim}

\subsection{Creazione della Pull Request}

\subsubsection{Step 1: Accedere a GitHub}
Navigare al repository su \url{https://github.com/eghosa02/mirofish_trading_offline}

\subsubsection{Step 2: Creare la Pull Request}
Dopo aver eseguito il push, GitHub mostrerà un pulsante \textbf{Compare \& pull request}. Cliccare su di esso.

In alternativa:
\begin{enumerate}
  \item Andare su \textbf{Pull requests}
  \item Cliccare su \textbf{New pull request}
  \item Selezionare il branch \texttt{feature\#<numero\_issue>} da confrontare con \texttt{main}
  \item Cliccare su \textbf{Create pull request}
\end{enumerate}

\subsubsection{Step 3: Compilare i dettagli della Pull Request}
\begin{itemize}
  \item \textbf{Titolo}: Deve essere descrittivo e conciso
  \item \textbf{Descrizione}: Includere:
  \begin{itemize}
    \item Riferimento alla issue (\#<numero>)
    \item Descrizione dei cambiamenti effettuati
    \item Eventuali note importanti
  \end{itemize}
\end{itemize}

Esempio di descrizione:
\begin{verbatim}
Closes #42

## Descrizione
Implementazione della nuova funzione di autenticazione.

## Cambiamenti
- Aggiunto modulo di login
- Aggiunta validazione credenziali
- Aggiunto test unitari

## Note
Testato su Windows 10
\end{verbatim}

\subsubsection{Step 4: Sottomissione e review}
\begin{enumerate}
  \item Cliccare su \textbf{Create pull request}
  \item Attendere la review del team
  \item Effettuare i cambiamenti richiesti se necessario
  \item Una volta approvata, cliccare su \textbf{Merge pull request}
\end{enumerate}

\subsection{Sincronizzazione e aggiornamenti}

Per mantenere il branch locale aggiornato con i cambiamenti del main:

\begin{verbatim}
git checkout main
git pull origin main
git checkout feature#<numero_issue>
git merge main
\end{verbatim}

Per risolvere conflitti (se presenti):
\begin{enumerate}
  \item Aprire i file con conflitti nell'editor
  \item Risolvere manualmente i conflitti
  \item Effettuare il commit dei cambiamenti
  \item Eseguire il push
\end{enumerate}
